% !TEX TS-program = pdflatexmk
\documentclass[12pt]{amsart}

%\usepackage[parfill]{parskip}    % Activate to begin paragraphs with an empty line rather than an indent

\usepackage[margin=1in]{geometry}

\usepackage{amsmath,amssymb,amsthm,latexsym,graphicx}
\usepackage[normalem]{ulem}
\usepackage{setspace} %used for doublespacing, etc.
\usepackage{hyperref}
\usepackage[shortlabels]{enumitem}
\usepackage{cancel}
\usepackage[dvipsnames,usenames]{color}
\usepackage[all]{xy}
\DeclareGraphicsRule{.tif}{png}{.png}{`convert #1 `dirname #1`/`basename #1 .tif`.png}

%Some useful environments.
\newtheorem{theorem}{Theorem}
\newtheorem{corollary}[theorem]{Corollary}
\newtheorem{conjecture}[theorem]{Conjecture}
\newtheorem{lemma}[theorem]{Lemma}
\newtheorem{proposition}[theorem]{Proposition}
\newtheorem{definition}[theorem]{Definition}
\newtheorem{example}[theorem]{Example}
\newtheorem{axiom}{Axiom}
\theoremstyle{remark}
\newtheorem{remark}{Remark}
\newtheorem*{exercise}{Exercise}%[section]

%Some useful shortcuts for our favorite fields
\newcommand{\RR}{{\mathbb R}}
\newcommand{\NN}{{\mathbb N}}
\newcommand{\ZZ}{{\mathbb Z}}
\newcommand{\QQ}{{\mathbb Q}}
\newcommand{\CC}{{\mathbb C}}

%Some useful shortcuts for formatting lists
\newcommand{\bc}{\begin{center}}
\newcommand{\ec}{\end{center}}
\newcommand{\be}{\begin{enumerate}}
\newcommand{\ee}{\end{enumerate}}
\newcommand{\bi}{\begin{itemize}}
\newcommand{\ei}{\end{itemize}}

%Some useful shortcuts for formatting mathematical symbols
\newcommand{\ol}[1]{\overline{#1}}
\newcommand{\oimp}[1]{\overset{#1}{\Longleftrightarrow}} %labeled iff symbol
\newcommand{\bv}[1]{\ensuremath{ \vec{\mathbf{#1}}} } %makes a vector.
\newcommand{\mc}[1]{\ensuremath{\mathcal{#1}}} %put something in caligraphic font
\newcommand{\mpg}[1]{\marginpar{ #1}} %to put comments in margins
\newcommand{\bsl}[1]{\texttt{\symbol{92}{\em #1}}} %for backslashes.
\newcommand{\normale}{\trianglelefteq}
\newcommand{\normal}{\triangleleft}

%Commenting tools
\usepackage{soul}
\definecolor{highlight}{rgb}{1,0.6,0.6}
\sethlcolor{highlight}
\newcommand{\hlm}[1]{\colorbox{highlight}{$\displaystyle #1$}}
\newtheoremstyle{mycomment}{\topsep}{-0in}{\small \itshape \sffamily}{}{\small \itshape\sffamily}{:}{.5em}{}
\theoremstyle{mycomment}
\newtheorem*{acomment}{\color{BrickRed}{Comment}}
\newcommand{\com}[1]{{\color{OliveGreen}\begin{acomment}{#1} %#2 \color{black} 
\end{acomment}\noindent}}
%\newcommand{\com}[1]{{\color{BrickRed}{\\ Comment:}\color{OliveGreen}{#1} \\}}
\newcommand{\red}[1]{{\color{BrickRed} #1}}
\def\midd{\,\middle\vert\, }
\newcommand{\blue}[1]{{\color{MidnightBlue}#1}}
\newcommand{\green}[1]{{\color{OliveGreen}#1}}
\newcommand{\mwrong}[2]{\red{\cancel{#1}}\green{#2}}
\newcommand{\wrong}[2]{\red{\sout{#1}}\green{#2}}
\definecolor{OliveGreen}{rgb}{.3,.5,.2}
\definecolor{MidnightBlue}{rgb}{.3,.4,.6}

\title{Assignment \#11 \today}

\author{Coordinator: Sasha}
\begin{document}
\maketitle
%\setstretch{2.5} %Use for 2.5 spacing
\doublespacing %Use for double spacing


\section*{Section 7.4}

\begin{exercise}[\textbf{7.4.7} --- Stefano Fochesatto] Let $R$ be a commutative ring with 1. Prove that the principal ideal generated by $x$
  in the polynomial ring $R[x]$ is a prime ideal if and only if $R$ is an integral domain. Prove that $(x)$ is a 
  maximal ideal if and only if $R$ is a field. 
  \begin{proof} Let $R$ be a commutative ring with 1. Consider the homomorphism $\varphi: R[x] \to R$ defined by 
    $\varphi(a_0 + a_1x + a_2x^2 + \dots) = a_0$. Let $p(x), q(x) \in R[x]$ such that $\varphi(p(x)) = a_0$ and $\varphi(q(x)) = b_0$.
    Note that by polynomial multiplication $\varphi(p(x)q(x)) = a_0b_0 = \varphi(p(x))\varphi(q(x))$ and clearly $\varphi$ is a surjection.
    Also note that every $p(x) \in \ker(\phi)$ must have zero constant term and therefore is comprised of products and sums of $x$ so $\ker(\varphi) = (x)$.

    By The First Isomorphism Theorem for Rings we find that $R[x]/(x) \cong R$. By Proposition 13, supposing $(x)$ is a prime ideal in $R[x]$ we know that $R[x]/(x)$ is an integral domain, thus $R$ is an integral domain. 
    Since Proposition 13 is an if and only if statement if we suppose $R$ is an integral domain then we get that $R/(x)$ is an integral domain and therefore $(x)$ is a prime ideal. 

    Supposing $(x)$ is a maximal ideal we get that $R/(x)$ is a field by Proposition 12 and therefore $R$ is a field. 
    Again since Proposition 12 is an if and only if statement, supposing $R$ is a field we know that $R/(x)$ is a field and therefore 
    $(x)$ must be a maximal ideal. 

  \end{proof}
\end{exercise}
\vspace{.5in}


\begin{exercise}[\textbf{7.4.41} --- Stefano Fochesatto] A proper ideal $Q$ of the commutative ring $R$ is called primary if whenever 
  $ab \in Q$ and $a \not\in Q$ then $b^n  \in Q$ for some positive integer $n$. Establish the following facts about 
  primary ideals, 
  \begin{enumerate}
  \item[(a)] The primary ideals of $\ZZ$ are $0$ and $(p^n)$, where $p$ is a prime and $n$ is a positive integer.
  \begin{proof} Clearly $0$ is a primary ideal of $\ZZ$. Let $ab \in (p^n)$ then either $p\mid a$ or $p\mid b$ which implies either $a^n  \in (p^n)$ or $b^n  \in (p^n)$.
    Suppose for the sake of contradiction $(i)$ is a prime ideal where $i$ is not prime or prime power. Let $i = p_1^{\alpha_1}p_2^{\alpha_2} \dots p_n^{\alpha_n}$ where $p_k$ are prime and 
    note that by definition $(p_1^{\alpha_1}p_2^{\alpha_2} \dots p_{n - 1}^{\alpha_{n - 1}})(p_n^{\alpha_n}) = i \in (i)$ however clearly neither $p_n^{\alpha_n} \in (i)$ or $(p_1^{\alpha_1}p_2^{\alpha_2} \dots p_{n - 1}^{\alpha_{n - 1}}) \in (i)$
    so $(i)$ is not primary. 
  \end{proof}
\vspace{.15in}

\item[(b)] Every prime ideal of $R$ is a primary ideal. 
\begin{proof} Suppose $P$ is a prime ideal. By definition if $ab \in P$ then either $a \in P$ or $b \in P$. Without loss of generality suppose 
  $a \not\in P$ and $b \in P$, then clearly $b^n \in P$ and we find that $P$ is primary. 
\end{proof}
\vspace{.15in}

\item[(c)] An ideal $Q$ of $R$ is primary if and only if every zero divisor in $R/Q$ is a nilpotent element of $R/Q$. 
\begin{proof} Suppose $Q$, an ideal of $R$ is primary. Let $(i + Q) \in R/Q$ be a zero divisor, so $(j + Q) \in R/Q$ such that $(i + Q)(j + Q) = (ij + Q) = Q$ where $(i + Q),(j + Q) \neq Q$. Therefore 
  $ij \in Q$ and $j \not\in Q$ since $Q$ is primary, there exists some $i^n \in Q$ so we get $(i + Q)^n = (i^n + Q) = Q$. Thus $(i + Q)$ is nilpotent in $R/Q$. 

  Suppose every zero divisor in $R/Q$ is a nilpotent element of $R/Q$. Let $ab \in Q$ and $a \not\in Q$ note that $(ab + Q) = (a + Q)(b + Q) = (a + Q)(Q) = Q$. Thus 
  $(b + Q)$ is a zero divisor. Then there exists some $(b + Q)^n = (b^n + Q) = Q$ so $b^n \in Q$. Therefore $Q$ is primary. 
\end{proof}
\vspace{.15in}

\item[(d)] If $Q$ is a primary ideal then $\text{rad}(Q)$ is a prime ideal. 
\begin{proof} Suppose $Q$ is a primary ideal. Let $ab \in \text{rad}(Q)$. By definition $\text{rad}(Q)$ there exists some $n$ such that $(ab)^n \in Q$, 
  Since $R$ is commutative we know that $(a^n)(b^n) \in Q$, without loss of generality suppose $(a^n) \not\in Q$, since $Q$ is primary then there exists some 
  $b^{mn} \in Q$. Therefore $b \in \text{rad}(Q)$.
\end{proof}


  \end{enumerate}
\end{exercise}
\vspace{.5in}





\end{document}